%% Appendices (shared).
\section{The computations}\label{app:computation}

This appendix describes the computation behind Theorems~\ref{thm:small}
and~\ref{thm:unique}. All
arithmetic is exact integer arithmetic; no floating-point comparison enters
any of the reported checks.

\subsection*{The recurrence}
Let $T$ be a tree rooted at a vertex $\rho$, and for a vertex $v$ let $T_v$
be the subtree induced by $v$ and its descendants. Consider the subsets $S$
of $V(T_v)$ such that every vertex of $T_v$ other than $v$ is in $S$ or has a
neighbour in $S$. Let $A_v$, $B_v$ and $C_v$ be the generating polynomials, by
cardinality, of those among these sets for which, respectively, $v\in S$;
$v\notin S$ and some child of $v$ is in $S$; $v\notin S$ and no child of $v$
is in $S$. If $v$ is a leaf, then $(A_v,B_v,C_v)=(x,0,1)$, and otherwise,
with $w$ running over the children of $v$,
\begin{equation}\label{eq:dp}
  A_v=x\prod_w(A_w+B_w+C_w),\qquad
  B_v=\prod_w(A_w+B_w)-\prod_w B_w,\qquad
  C_v=\prod_w B_w .
\end{equation}
Indeed, if $v\in S$ the children are unconstrained; if $v\notin S$, every
child must be dominated inside its own subtree, and $C_v$ collects the sets in
which no child is chosen. The root must be dominated, so
$D(T,x)=A_\rho+B_\rho$.

\subsection*{Two programs}
The computation was performed twice, by programs written independently of
each other.

\emph{Program I} generates all rooted trees of order $n$ as level sequences
by the algorithm of Beyer and Hedetniemi~\cite{BH}, and checks that their
number is the number $r_n$ of rooted trees given by the classical recurrence
$(n-1)r_n=\sum_{j=1}^{n-1}\bigl(\sum_{d\mid j}d\,r_d\bigr)r_{n-j}$. It
keeps the rooted trees whose root is a centroid; when the tree has two
centroids, the two rootings are compared through a canonical code and one of
them is kept. Every tree is therefore processed exactly once, and the program
checks that the number of trees processed equals the number
\[
  t_n=r_n-\tfrac12\Bigl(\sum_{k=1}^{n-1}r_kr_{n-k}-[\,n\text{ even}\,]\,r_{n/2}\Bigr)
\]
of trees of order $n$ given by Otter's formula~\cite{Otter}. The domination
polynomial is computed by merging the children of a vertex one at a time,
which is~\eqref{eq:dp} in the equivalent form
$(A,B,C)\mapsto\bigl(A(A_w+B_w+C_w),\,B(A_w+B_w)+CA_w,\,CB_w\bigr)$ starting
from $(x,0,1)$. For every tree of order at most $16$ the program also
enumerates all vertex subsets and compares the resulting numbers of
dominating sets and the values of $\gamma$ and $\Gamma$ with those obtained
from the recurrences (the program computes $\Gamma$ as the independence
number, cf.~\eqref{eq:alpha-Gamma}), and it checks the two
inequalities~\eqref{eq:twosided} for every dominating set; there are $32\,507$ such trees and no
discrepancy.

\emph{Program II} generates the trees of order $n$ directly, by the
algorithm of Wright, Richmond, Odlyzko and McKay~\cite{WROM}, checks their
number against $t_n$, and evaluates~\eqref{eq:dp} as written.

Both programs test $d_j^2\ge d_{j-1}d_{j+1}$ for all $1\le j\le n-1$ with
128-bit integer products; the coefficients themselves are smaller than
$2^{27}$ in the range considered. For each order, both programs also output the
sum, modulo $2^{64}$, of a 64-bit hash of the coefficient vector over all
trees of that order. These sums agree for every order from $2$ to $\MaxOrder$,
which shows, up to hash collisions, that the two programs computed the same
multiset of polynomials although they traverse the trees in different orders
and with different rootings.

\subsection*{Results}
Table~\ref{tab:counts} lists the numbers of trees. Together with the tree on
one vertex, whose domination polynomial is $x$, these are the $\TotalTreesAll$
trees with at most $\MaxOrder$ vertices. Both programs report
that the domination polynomial is log-concave for every tree with at most
$\MaxLC$ vertices, and that exactly one tree with $\MinNonLC$ vertices, the
tree $U$ of Figure~\ref{fig:U}, fails to have this property. Both programs
also report that all these domination polynomials are unimodal.

\begin{table}[ht]
\centering
\caption{Numbers $r_n$ of rooted trees generated by Program~I and numbers
$t_n$ of trees examined by both programs.}\label{tab:counts}
\smallskip
\begin{tabular}{rrr@{\qquad\qquad}rrr}
\hline
$n$ & $r_n$ & $t_n$ & $n$ & $r_n$ & $t_n$\\
\hline
2 & 1 & 1 & 15 & 87\,811 & 7\,741\\
3 & 2 & 1 & 16 & 235\,381 & 19\,320\\
4 & 4 & 2 & 17 & 634\,847 & 48\,629\\
5 & 9 & 3 & 18 & 1\,721\,159 & 123\,867\\
6 & 20 & 6 & 19 & 4\,688\,676 & 317\,955\\
7 & 48 & 11 & 20 & 12\,826\,228 & 823\,065\\
8 & 115 & 23 & 21 & 35\,221\,832 & 2\,144\,505\\
9 & 286 & 47 & 22 & 97\,055\,181 & 5\,623\,756\\
10 & 719 & 106 & 23 & 268\,282\,855 & 14\,828\,074\\
11 & 1\,842 & 235 & 24 & 743\,724\,984 & 39\,299\,897\\
12 & 4\,766 & 551 & 25 & 2\,067\,174\,645 & 104\,636\,890\\
13 & 12\,486 & 1\,301 & 26 & 5\,759\,636\,510 & 279\,793\,450\\
14 & 32\,973 & 3\,159 & 27 & 16\,083\,734\,329 & 751\,065\,460\\
\hline
\end{tabular}

\end{table}

The source code of both programs, their complete outputs, a script that
recomputes the number in Remark~\ref{rem:boundary}, and a program that
enumerates all subsets of a given tree (used for
Proposition~\ref{prop:U}) are contained in the supplementary material.

\section{Formal verification}\label{app:lean}

The following statements have been formalised and checked in
Lean~4 (version 4.27.0) with Mathlib~\cite{Lean,Mathlib}. In the formal
development a graph is a term of Mathlib's type \texttt{SimpleGraph} on a
finite vertex type, a forest is an acyclic graph, and the numbers $\gamma$
and $\Gamma$ are specified by explicit witnesses and bounds rather than by a
defined invariant; for example, $\gamma$ is a natural number such that some
dominating set has at most $\gamma$ elements and every dominating set has at
least $\gamma$ elements.

\begin{itemize}
\item Lemma~\ref{lem:lower}
(\lean{forest_critical_inequality}) and Lemma~\ref{lem:upper}
(\lean{bipartite_critical_le_two_independence}).
\item The core of Theorem~\ref{thm:main}
(\lean{graph_uniform_gap_unimodal}): for an arbitrary relation on a
finite type, natural numbers $h<K$ with $(h+1)^2\le4K$, and an index $A$ with
$\gamma+K\le A+2\le n$ up to which the coefficients are nondecreasing,
condition $(\mathrm{H}_h)$ implies unimodality.
\item Corollary~\ref{cor:gapthree} for forests without isolated vertices and
gap exactly three (\lean{forest_gap_three_unimodal_no_isolates}).
\item Theorem~\ref{thm:criterion-intro}(i) for $h=4$
(\lean{forest_gap_four_large_unimodal}), and the fact that a graph
without isolated vertices that has a minimal dominating set of size
$\gamma+4$ satisfies $n\ge\gamma+14$ or $n\le22$
(\lean{no_isolates_gap_four_large_or_at_most_twenty_two}).
\end{itemize}

All these theorems depend only on the three standard axioms of Lean's
classical logic (\lean{propext}, \texttt{Classical.choice},
\texttt{Quot.sound}). The formalisation does \emph{not} cover the reduction to
forests without isolated vertices (Lemma~\ref{lem:isolates}), the statement
of Theorem~\ref{thm:criterion-intro} for general $h$ in the form given here,
Lemma~\ref{lem:product}, or any part of the computation of
Appendix~\ref{app:computation}. In particular, Theorems~\ref{thm:small} and~\ref{thm:unique}, and
with the former the case of gap four of Theorem~\ref{thm:gapfour-intro}, are
computer-assisted results in the usual sense and not formally verified ones.
The Lean sources are part of the supplementary material.

